\documentclass[10pt]{article}

\usepackage[utf8]{inputenc}

\usepackage[T1]{fontenc}

\usepackage{graphicx}

\usepackage[export]{adjustbox}

\graphicspath{ {./images/} }

\usepackage{amsmath}

\usepackage{amsfonts}

\usepackage{amssymb}

\usepackage[version=4]{mhchem}

\usepackage{stmaryrd}

\usepackage{bbold}



\begin{document}

ЦЕПНАЯ ЛИНИЯ - плоская трансцендентная кривая, форму к-рой принимает под действием силы тя-



\includegraphics[max width=\textwidth, center]{2023_07_09_c21558d40729e740c28eg-1(1)}

жести однородная гибкая нерастяжимая тяжелая нить с закрепленными концами (см. рис.). Уравнение в декартовых прямоугольных координатах:



\[

y=\frac{a}{2}\left(e^{x / a}+e^{-x / a}\right)=a \operatorname{ch} \frac{x}{a} .

\]



Длина дуги, отсчитываемая от точки $x=0$ :



\[

l=\frac{1}{2}\left(e^{x / a}-e^{-x / a}\right)=a \operatorname{sh} \frac{x}{a} .

\]



Радиус кривизны:



\[

r=a \operatorname{ch}^{2} \frac{x}{a}

\]



Площадь, ограничиваемая дугой Ц. л., двумя её ординатами и осью абсцисс:



\begin{center}

\includegraphics[max width=\textwidth]{2023_07_09_c21558d40729e740c28eg-1(3)}

\end{center}



Если дугу Ц. л. вращать вокруг оси абсцисс, то образуется катеноид.



Лит.: [1] С а в е л о в А. А., Плоские кривые, М., 1960.



ЦЕПНАЯ РЕКУРРЕНТНОСТЬ - наиболе из свойств «Повторяемости движений», рассматриваемое в топологической динамике. В основном случае топологич. потока $\left\{S_{t}\right\}$ на метрич. компакте $W$ с метрикой $\rho$ точка $w \in W$ обладает свойством Ц. р., если для любых $\varepsilon, T>0$ имеется $\varepsilon$-траектория, исходящая из $w$ и снова возвращающаяся в $w$ через время $T_{\varepsilon}>T$. Под $\varepsilon$-траекторией понимается такая параметризованная кривая (возможно, разрывная) $w(t), \quad 0 \leqslant t \leqslant T_{\varepsilon}$, что $\rho\left(S_{\tau} w(t), w(t+\tau)\right)<\varepsilon$ при $0 \leqslant \tau \leqslant 1,0 \leqslant t \leqslant T_{\varepsilon}-1$ («конечный отрезок $\varepsilon$-траектории близок к отрезку настоящей траектории»). Имеется определение Ц. р. и для более общего случая [1]. Если $W$ является замкнутым многообразием, то Ц. p. совпадает со свойством «слабой неблуждаемости» (см. [3]), более непосредственно отражающим влияние малых (в топологич. смысле) возмущений системы на поведение ее траекторий. Вне множества точек со свойством Ц. р. поведение системы напоминает поведение градиентной динамической системь (cm. [1], [2]).



Jum.: [1] C o n l e y C h., Isolated invariant sets and the Morse index, Providence, 1978; [2] S h u b M., "Astérisque",



\includegraphics[max width=\textwidth, center]{2023_07_09_c21558d40729e740c28eg-1}

с к и й В. А., в кн.: Динамические системы и вопросы устойчи$\begin{array}{lll}\text { вости решений дифференциальных уравнений, К., } 1973, \\ \text { с. } 165-74 . & \text { Д. В. Аиосов. }\end{array}$



ЦЕПНОЕ КОЛЬЦО л е в о е - колыцо (обыно предполагаемое ассоциативным и с единицей), левые идеалы к-рого образуют цепь. Другими словами, $R$ левое Ц. к., если $R$ - левый цепной модуль нал собоӥ. Всякое левое Ц. к. локально. Аналогично определя отся правые Ц. к. Коммутативные Ц. к. часто называют кольц а и и о р и рова ни я. См. также Дискретного норлирования кольцо. Л. А. Скорияков.



ЦЕПНОИ МОДУЛЬ - модуль, совокупность всех подмодулей к-рого образует цегь, т. є. линейно упорядоченное множество; для этого достаточно, чтобы цепью была совокупность всех циклич. подмодулеї.



ЦЕПЬ - 1) то жке, что линейно упорядоченное множество.



\begin{enumerate}

  \setcounter{enumi}{1}

  \item Ц.- формальная линейная комбинация симплексов (триангуляции, симплициального множества и, в частности, сингулярных симплексов топологич. пространства) или клеток. В более общем смысле - элемент группы цепей произвольного цепного комплекса (как правило, свободного). If . с коэффициентами в грушпе $G$ - әлемент тензорного произведения депного комплекса на группу $G$. Лum.: [1] С т ин р о д Н., Э й л е н б е р г С., Основания алгебраической топологии, пер. с англ., М., 1958; [2] X и Лт о н П.-Д ж., У а й л и С., Теория гомологий, пер. с англ., М.,

\end{enumerate}



1966.



ЦЕРМЕЈО АКСМОМА - выбора аксиома для произвольного (не обязательно дизыюнктного) семейства множеств. Эту аксиому Э. Цермело сформулировал в 1904 в виде следующего утверждения, названного им п р и ц И І о м в Ы 6 о р а [1]: для любого семейства множества $t$ можно выбрать из каждого его члена по единственному представителю и объединить их всех в одно множество, - и дал первое доказательство своей теоремы о вполне упорядочении. В 1906 Б. Рассел (В. Russell) сформулировал аксиому выбора в мультипликативной форме: если $t$ есть дизыюнктное множество непустых множеств, то прямое произведение ПI не пусто. В 1908 Э. Цермело доказал әквивалентность мультипликативной формы аксиомы выбора ее обычной формулировке.



Jum.: [1] Z e r m e lo E., "Math. Ann.», 1904, Bd 59, S. $514-16 ;[2]$ [ р е н к е л ь А., Б а p-Х и л л е л И., Основания теории множеств, пер. с англ., М., 1966.



\includegraphics[max width=\textwidth, center]{2023_07_09_c21558d40729e740c28eg-1(2)}

вполне упорядочить (см. Вполне упорядоченное множество). Впервые әту теорему доказал Ә. Цермело (Е. Zermelo, 1904), исходя из принципа выбора - одной из әквивалентных форм аксиомы выбора (см. Церлело аксиома). Позднее выяснилось, что Ц. т. эквивалентна аксиоме выбора (в системе обычных аксиом теории множкеств), а значит, и многим другим высказываниям теоретико-множественного характера (см. Bыбора аксиома).



ЦИКЈ - чеnь, граница к-рой равна нулю. A. Ф. Харициладзе.



ЦИКЛИЧЕСКАЯ ГРУППА - групша с одним образующим. Все Ц. г. абелевы. Всякая конечная группа простого порядка - Ц. г. Супествует по одной, с точностью до изоморфизма, Ц. г. каждого конечного порядка $n$ и одна бесконечная Ц. г., изоморфная аддитивной группе $\mathbb{Z}$ целых чисел. Конечная Ц. Г. $G$ порядка $n$ изоморфна аддитивной групте кольца вычетов $\mathbb{Z}(n)$ по модулю $n$ (а также группе $\mathbb{C}(n)$ (комплексных) корней степени п из 1). В качестве образующего в этой группе может быть взят любой элемент $a$ порядка $n$. Тогда



\[

G=\left\{1=a^{0}=a^{n}, a, a^{2}, \ldots, a^{n-1}\right\} .

\]



О. А. Иванова. ДИКЛИЧЕСКАЯ ПОЛУГРУППА - то не, что моногенная полугруппа.



ЦИКЛИЧЕСКИЕ КООРДИНАТБІ-обобщенные координаты нек-роӥ физич. спстемы, не входящие явно в выражение характеристич. функции этой системы. При использовании соотвтствующих уравнений движения можно сразу получить для каждой Ц. к. соответствующий ей интеграл движения. Напр., если Лагранжа gi ункция $L,\left(q_{i}, \dot{q_{i}}, t\right)$, где $q_{i}$ - обобщенные координаты, $q_{i}$ - обобщенные скорости, $t$ - время, не содержит явно $q_{j}$, то $q_{j}-$ Ц.к. и j-е Лагранжа уравнение имеет вІІ $\frac{d}{d t}\left(\frac{\partial L}{\partial \dot{q}_{j}}\right)=0$, что сразу дает интеграл двІжения



\[

\frac{\partial L}{\partial \dot{q}_{j}}=\text { const. }

\]



ЦИКЛИЧЕСКИЙ МОДУЛЬ н а д ко Л. Ц. Сололов. (л е в ы й) - фактормодуль кольца $R$, рассматриваемого как левый $R$-модуль, по нек-рому левому идеалу. В частности, циклическими являются неприводимые модули. С Ц. м. связана п $\mathrm{p}$ о б л е м а $\mathrm{K}$ ё т е (см. [4]): над какими колыцами каждый (или каждый конеqно порожденный) модуль изоморфен прямой сумме Ц. м.? В классе коммутативных колец таковыми ока-





\end{document}